\chapter{Louis Nirenberg (2015)}\index{Nirenberg, Louis}

\section{Biographical Sketch}

Louis Nirenberg was born on 28 February 1925 in Hamilton, Ontario, Canada. He studied at McGill University and New York University, where he completed his PhD in 1949 under the supervision of James Serrin and Kurt Friedrichs. He spent his entire career at the Courant Institute of Mathematical Sciences at NYU. Nirenberg received the Abel Prize (jointly with John F. Nash Jr.) in 2015, the National Medal of Science (1995), and the Chern Medal (2010).

\section{High-Level View for a General Audience}

\subsection{What Did Nirenberg Do?}

Think about the shape of a soap film stretched across a wire frame. The film takes the shape that minimizes its surface area. This is a problem of geometry: given a boundary, find the surface of least area spanning it.

Now imagine that you are told the curvature of a surface at every point, but not its overall shape. Can you reconstruct the shape from this curvature information? This is the kind of problem that Louis Nirenberg solved.

Nirenberg made fundamental contributions to the theory of partial differential equations\index{PDEs}, particularly equations describing surfaces and their curvatures. He developed the tools needed to prove the existence and regularity of solutions to these geometric PDEs\index{PDEs}. His estimates and methods---including Schauder estimates, the Agmon--Douglis--Nirenberg (ADN) theory, the Gagliardo--Nirenberg inequalities, and the method of moving planes---are among the most widely used tools in analysis.

\subsection{Why Does It Matter?}

\begin{itemize}
\item \textbf{Geometry and Physics}: Nirenberg's PDE estimates are the workhorses behind the existence and regularity of solutions to the geometric equations that describe minimal surfaces, Ricci flow, and the shape of the universe itself.

\item \textbf{Materials Science}: The mathematics of surfaces and their curvatures, which Nirenberg studied, is used to understand the behavior of thin films, crystals, and other materials.

\item \textbf{Image Processing}: The partial differential equations that Nirenberg studied are used in image processing, including edge detection, denoising, and image segmentation.

\item \textbf{Animation and Graphics}: The reconstruction of surfaces from geometric data, which uses the mathematics of Nirenberg, is used in computer graphics and animation.
\end{itemize}

\section{The Origin of the Problem}

\subsection{The Geometric PDE Problem}

Many problems in differential geometry reduce to solving partial differential equations. For example, the prescribed curvature problem asks: given a function $K$ on a surface, can we find a metric whose Gaussian curvature is $K$? This leads to the Yamabe equation\index{Yamabe equation} and the Monge--Amp\`ere equation\index{Monge--Amp\`ere equation}.

These equations are \emph{nonlinear} and \emph{elliptic}. Their analysis requires powerful techniques that Nirenberg helped develop.

\begin{primer}{What Are Elliptic PDEs?}
Elliptic partial differential equations are the mathematical counterpart to steady-state physical problems. Laplace's equation $\nabla^2 u = 0$, for example, describes the equilibrium temperature distribution on a metal plate. What makes elliptic equations special is that their solutions are automatically smooth---they have no corners or jumps---and are determined entirely by boundary data. Nirenberg's estimates provide precise control over how smooth the solutions are, which is essential for proving that geometric surfaces with prescribed curvature actually exist.
\index{Regularity theory}
\end{primer}

\section{Deep Insight into the Problem}

\subsection{Nirenberg's Work on Elliptic PDEs}

Nirenberg made fundamental contributions to the theory of linear and nonlinear elliptic PDEs:
\begin{itemize}
\item \textbf{Schauder Estimates}: Nirenberg developed a priori estimates for solutions of elliptic equations of arbitrary order, including the interior Schauder estimates and the boundary estimates. These estimates control the $C^{k,\alpha}$ norm of the solution in terms of the $C^{k-2,\alpha}$ norm of the equation.
\item \textbf{Agmon--Douglis--Nirenberg (ADN) Theory}: With Agmon and Douglis, Nirenberg developed a general theory of elliptic boundary value problems for systems of equations. The ADN theory is the standard framework for analyzing elliptic boundary value problems.
\item \textbf{The Method of Moving Planes}: Nirenberg developed this technique for proving symmetry of solutions to nonlinear PDEs. The idea: reflect the solution across a plane and compare the original with the reflection using the maximum principle\index{Maximum principle}. If the solution is monotone up to a critical plane, it must be symmetric.
\end{itemize}

\subsection{The Gagliardo--Nirenberg--Sobolev Inequality}

This fundamental inequality controls the $L^{p^*}$ norm of a function in terms of its gradient:
\[
\|u\|_{L^{p^*}(\mathbb{R}^n)} \leq C \|\nabla u\|_{L^p(\mathbb{R}^n)},
\]
where $p^* = \frac{np}{n-p}$ is the Sobolev conjugate exponent. This inequality is essential for establishing compactness in PDE theory.

\section{Biggest Contributions and New Tools}

\subsection{Nirenberg's Maximum Principle and the Moving Planes Method}

Nirenberg's maximum principle for elliptic equations states that the maximum of a solution occurs on the boundary unless the solution is constant. This simple but powerful tool is used throughout PDE theory.

\subsection{The Gagliardo--Nirenberg Interpolation Inequality}

This inequality estimates intermediate derivatives in terms of lower and higher derivatives:
\[
\|D^j u\|_{L^p} \leq C \|D^m u\|_{L^r}^\theta \|u\|_{L^q}^{1-\theta}.
\]

\subsection{The Brezis--Nirenberg Problem}

The Brezis--Nirenberg problem concerns the existence of positive solutions to the critical Sobolev equation:
\[
-\Delta u = u^{(n+2)/(n-2)} + \lambda u \quad \text{in } \Omega, \quad u = 0 \text{ on } \partial\Omega.
\]

\section{Impact of the Discovery}

\subsection{Impact on Mathematics}

\begin{itemize}
\item \textbf{PDE Theory}: Nirenberg's work on elliptic PDEs provides essential tools for geometric analysis, the theory of minimal surfaces, and the study of complex Monge--Amp\`ere equations.

\item \textbf{Functional Analysis}: Nirenberg's inequalities and estimates are foundational to modern analysis.

\end{itemize}

\begin{whyitmatters}
Nirenberg's PDE estimates are foundational: they are the workhorses behind the existence and regularity of solutions to the geometric equations that describe minimal surfaces, Ricci flow, and the shape of the universe itself. Together with his coauthors, he built the framework that makes modern geometric analysis possible.
\end{whyitmatters}

\subsection{Impact on Technology}

\begin{itemize}
\item \textbf{Computer Vision}: The level set method and the Perona--Malik equation, based on PDEs of the type studied by Nirenberg, are used in image processing.

\item \textbf{Materials Science}: The geometry of thin films and interfaces uses the theory of surfaces and curvature that Nirenberg advanced.
\end{itemize}

\section{Current Developments}

\subsection{Nonlinear PDEs in Geometry}

Nirenberg's methods continue to be essential in the study of:
\begin{itemize}
\item \textbf{The Yamabe Problem}: Finding a metric of constant scalar curvature.
\item \textbf{The Prescribed Scalar Curvature Problem}: Given a function $K$, does there exist a metric with scalar curvature $K$?
\item \textbf{The Ricci Flow}: The parabolic version of the Yamabe problem, crucial for Perelman's proof of the Poincar\'e conjecture.
\end{itemize}

\subsection{The Brezis--Nirenberg Problem Today}

The Brezis--Nirenberg problem continues to be an active area of research. It was the starting point for the study of critical elliptic equations, and variations of the problem appear in the study of the Yamabe equation, the prescribing scalar curvature problem, and the study of elliptic equations on compact manifolds.

\section{Conclusion}

Louis Nirenberg transformed our understanding of partial differential equations. His work on elliptic PDEs provided the tools needed to solve the geometric equations that arise in the study of surfaces and their curvatures. Together with John F. Nash Jr., he received the Abel Prize ``for striking and seminal contributions to the theory of nonlinear partial differential equations and its applications to geometric analysis.''

\section{Behind the Breakthrough: The Human Story}
Nirenberg shared the 2015 Abel Prize with John Nash, recognizing their parallel and complementary contributions to nonlinear PDEs and geometric analysis. Nirenberg's career at the Courant Institute spanned more than six decades, during which he advised generations of students and shaped the modern field of PDE theory.

\section{Pedagogical Metaphors and Lineage}
\subsection{Signature Metaphor}
``Building a scaffold of estimates that turns a slippery PDE problem into a well-posed, solvable system---like adding handrails to a steep mountain trail.''

\subsection{Mathematical Lineage}
Nirenberg advised by James Serrin and Kurt Friedrichs.

\section{Applied Science and Computational Algorithms}
\subsection{Key Takeaways for Applied Science}
Applied in computer graphics (isometric texture mapping, mesh deformation, and skinning) and in aerospace engineering to solve fluid flow and heat dissipation boundaries.

\subsection{Algorithms and Numerical Implementations}
Finite element methods for solving elliptic boundary value problems rely on Nirenberg-type estimates for convergence analysis and error bounds.

\section{The Research Frontier}
Proving global regularity or blow-up of solutions to the 3D Navier-Stokes equations, which is a major open problem in mathematical physics.

\section{Guided Exercises and Challenges}
\begin{enumerate}[label=\textbf{\arabic*.}]
  \item State the Gagliardo-Nirenberg-Sobolev inequality and explain its role in establishing compactness in PDE theory.
  \item Explain the method of moving planes and how the maximum principle is used to prove symmetry of solutions.
  \item Describe the Brezis-Nirenberg problem and its significance for critical Sobolev exponents.
\end{enumerate}

\section{Curriculum Map and Further Reading}
For students wishing to delve deeper into the mathematical machinery underlying these breakthroughs, the following learning path is recommended:
\begin{itemize}
  \item L. Nirenberg, \emph{Topics in Nonlinear Functional Analysis}, AMS.
  \item S. Agmon, A. Douglis, and L. Nirenberg, ``Estimates near the boundary for solutions of elliptic partial differential equations,'' CPAM 12 (1959).
\end{itemize}

\section*{Bibliography}
\begin{enumerate}
\item L. Nirenberg, ``Topics in Nonlinear Functional Analysis,'' New York University, 1974.
\item S. Agmon, A. Douglis, and L. Nirenberg, ``Estimates near the boundary for solutions of elliptic partial differential equations,'' Communications on Pure and Applied Mathematics 12 (1959), 623--727.
\item L. Nirenberg, ``The Weyl and Minkowski problems in differential geometry,'' Communications on Pure and Applied Mathematics 6 (1953), 337--394.
\item L. Nirenberg, ``On elliptic partial differential equations,'' Annali della Scuola Normale Superiore di Pisa 13 (1959), 115--162.
\item H. Brezis and L. Nirenberg, ``Positive solutions of nonlinear elliptic equations involving critical Sobolev exponents,'' Communications on Pure and Applied Mathematics 36 (1983), 437--477.
\item J. F. Nash Jr., ``The imbedding problem for Riemannian manifolds,'' Annals of Mathematics 63 (1956), 20--63.
\item R. S. Hamilton, ``The inverse function theorem of Nash and Moser,'' Bulletin of the American Mathematical Society 7 (1982), 65--222.
\end{enumerate}
